% from ega4-21b-fr.tex
\oldpage[IV-4]{288}
que $W^{(x)}=V^{(x)}$, autrement dit que $V^{(x)}$ est fermé dans $X$.
On peut alors définir un diviseur $D^{(x)}$ sur $X$ par la condition que $D^{(x)}|U^{(x)}=D'^{(x)}$, et $D^{(x)}|(X-V^{(x)})=0$, ce qui a un sens puisque dans $U^{(x)}\cap(X-V^{(x)})$ la restriction de $g^{(x)}$ est une section inversible, donc les restrictions à cet ouvert des deux diviseurs considérés sur $U^{(x)}$ et $X-V^{(x)}$ sont égales.
Il suffit alors, pour répondre à la question, de prendre $D=\sum_{x\in T}D^{(x)}$, ce qui a un sens puisque $T$ est fini.

  \item Compte tenu du diagramme commutatif \sref{IV.21.4.2.1}, la conclusion résultera de (i) si l'on prouve que tout $\sh{O}_{X_0}$-Module inversible $\sh{L}_0$ est isomorphe à un $\sh{O}_{X_0}$-Module de la forme $\sh{O}_{X_0}(D_0)$ \sref{IV.21.2.8}, où $D_0$ est un diviseur sur $X_0$ dont le support ne rencontre pas $Z_0$.
Or, comme $U_0$ est schématiquement dense dans $X_0$ \sref{IV.20.2.13}, (iv), il suffit pour cela de prouver qu'il existe une section $t\in\Gamma(U_0,\sh{L}_0)$ telle que $t(z_0)\neq0$ aux points de $\operatorname{Ass}(\sh{O}_{X_0})\cup Z_0$ \sref{IV.3.1.9}.
On peut donc se borner au cas où $X_0=\Spec(A_0)$ est affine;
mais alors $\sh{L}_0$ est ample \sref[II]{II.5.1.4} et comme l'ensemble $\operatorname{Ass}(\sh{O}_{X_0})\cup Z_0$ est fini, la conclusion résulte de \sref[II]{II.4.5.4}.
\end{enumerate}
\end{proof}

\begin{corollary}[21.9.12]
\label{IV.21.9.12}
Soient $A$ un anneau local hensélien \sref{IV.18.5.8}, $S=\Spec(A)$, $s_0$ le point fermé de $S$, $f:X\to S$ un morphisme séparé de présentation finie, et supposons que l'ensemble $X_0=f^{-1}(s_0)$ soit de dimension $\leq1$.
Alors, pour tout sous-schéma fermé $X'_0$ de $X$, ayant $X_0$ pour ensemble sous-jacent et de présentation finie sur $S$, l'homomorphisme canonique \sref{IV.21.3.2.4} $\operatorname{Pic}(X)\to\operatorname{Pic}(X'_0)$ est surjectif.
\end{corollary}

\begin{proof}
Montrons d'abord qu'il suffit de prouver le corollaire lorsque l'anneau local hensélien $A$ est de plus noethérien.
On sait en effet \sref{IV.18.6.15} que l'on peut écrire $A=\varinjlim A_\lambda$, où les $A_\lambda$ sont des anneaux locaux noethériens et henséliens, les homomorphismes $A_\lambda\to A$ étant locaux;
il existe en outre un indice $\alpha$ et un morphisme séparé de présentation finie $f_\alpha:X_\alpha\to S_\alpha=\Spec(A_\alpha)$ tels que $X$ et $f$ soient déduits de $X_\alpha$ et $f_\alpha$ par le changement de base $S\to S_\alpha$ \sref{IV.8.10.5}, (v);
avec les notations habituelles \sref{IV.8.8.1}, les morphismes $f_\lambda:X_\lambda\to S_\lambda$ seront séparés pour $\lambda\geq\alpha$ et on aura $X=X_\lambda\times_{S_\lambda}S$.
En outre, on peut supposer que, si $s_{0\alpha}$ est l'unique point fermé de $S_\alpha$, il y a un sous-schéma fermé $X'_{0\alpha}$ de $X_\alpha$, ayant même ensemble sous-jacent que $f_\alpha^{-1}(s_{0\alpha})$, tel que $X'_0=X'_{0\alpha}\times_{S_\alpha}S$ \sref{IV.8.6.3};
pour $\lambda\geq\alpha$, posons $X'_{0\lambda}=X'_{0\alpha}\times_{S_\alpha}S_\lambda$;
on a $\dim(X'_{0\lambda})=\dim(X'_0)\leq1$ par transitivité des fibres et \sref{IV.4.1.4}.
Il s'agit de voir que si l'on a prouvé que l'homomorphisme $\operatorname{Pic}(X_\lambda)\to\operatorname{Pic}(X'_{0\lambda})$ est surjectif pour tout $\lambda\geq\alpha$, il en est de même de $\operatorname{Pic}(X)\to\operatorname{Pic}(X'_0)$.
On a évidemment le diagramme commutatif d'homomorphismes canoniques
\[
  \xymatrix{
    \operatorname{Pic}(X_\lambda) \ar[r] \ar[d] & \operatorname{Pic}(X) \ar[d]\\
    \operatorname{Pic}(X'_{0\lambda}) \ar[r] & \operatorname{Pic}(X'_0)
  }
\]
Pour tout $\sh{O}_{X'_0}$-Module inversible $\sh{L}'_0$, il existe un $\lambda\geq\alpha$ et un $\sh{O}_{X'_{0\lambda}}$-Module inversible $\sh{L}'_{0\lambda}$ tel que $\sh{L}'_0$ s'en déduise par changement de base $S\to S_\lambda$ \sref{IV.8.5.2} et \sref{IV.8.5.5};
il suffit de considérer un $\sh{O}_{X_\lambda}$-Module inversible $\sh{L}_\lambda$ tel que $\operatorname{cl}(\sh{L}'_{0\lambda})$ soit l'image de $\operatorname{cl}(\sh{L}_\lambda)$ dans $\operatorname{Pic}(X'_{0\lambda})$, puis de prendre le $\sh{O}_X$-Module inversible $\sh{L}$ déduit de $\sh{L}_\lambda$ par changement de base;
en vertu de la commutativité du diagramme précédent, $\operatorname{cl}(\sh{L}'_0)$ sera l'image de $\operatorname{cl}(\sh{L})$.

Supposons donc $A$ noethérien, donc aussi $X$, et vérifions que $X$ et $X'_0$ satisfont aux conditions de \sref{IV.21.9.11}, (ii).
On a par hypothèse $\dim(X'_0)\leq1$;
d'autre part, pour vérifier la condition 2\textsuperscript{o} de \sref{IV.21.9.11}, considérons un sous-préschéma réduit $Y$ de $X$ ayant une partie localement fermée $Y$ pour ensemble sous-jacent;
le morphisme $g:Y\to S$, restriction de $f$, étant quasi-fini en chacun des points $x_i$ de $Y\cap X_0$ ($1\leq i\leq n$), on peut appliquer \sref{IV.18.5.11}, c), et on voit que $Y$ est somme des sous-préschémas ouverts $Y_i=\Spec(\sh{O}_{Y,x_i})$ ($1\leq i\leq n$) et d'un préschéma $Y''$ tel que $Y''\cap X_0=\emp$, et en outre que les injections canoniques $j_i:Y_i\to X$ sont telles que $f\circ j_i$ soit un morphisme fini.
Comme $f$ est séparé, $j_i$ est aussi un morphisme fini, donc $Y_i$ est fermé dans $X$;
on répond donc à la question en prenant $Y'=\bigcup_{1\leq i\leq n}Y_i$.

Il reste à vérifier l'hypothèse supplémentaire de \sref{IV.21.9.11}, (ii).
Or, $X'_0$ étant une courbe séparée sur le corps $\boldsymbol{k}(s_0)$, est un $\boldsymbol{k}(s_0)$-schéma quasi-projectif (chap.~V)\footnote{Le lecteur vérifiera que \sref{IV.21.9.12} n'est pas utilisé pour prouver cette propriété au chap.~V.}, donc il existe un $\sh{O}_{X'_0}$-Module ample \sref[II]{II.5.3.1}, et cela achève la démonstration.
\end{proof}

\begin{remark}[21.9.13]
\label{IV.21.9.13}
Sous les conditions de \sref{IV.21.9.12}, supposant en outre $f$ propre, le morphisme $f$ est même \emph{projectif}:
en effet, si $\sh{L}_0$ est un $\sh{O}_{X'_0}$-Module ample, il existe, en vertu de \sref{IV.21.9.12}, un $\sh{O}_X$-Module inversible $\sh{L}$ dont l'image réciproque dans $X'_0$ est isomorphe à $\sh{L}_0$, donc est ample.
Puisque tout voisinage de $s_0$ dans $S$ est nécessairement $S$ tout entier, on déduit alors de \sref{IV.9.6.4} que $\sh{L}$ est un $\sh{O}_X$-Module ample, d'où la conclusion (\sref[II]{II.5.3.1} et \sref[II]{II.5.5.3}).
\end{remark}

\subsection{Images réciproques et images directes de cycles 1-codimensionnels}
\EGAReaderTarget{ega4.c04.subsec.21-10}{subsection:IV.21.10}

Dans un chapitre ultérieur, consacré à la théorie des intersections, seront développées de façon systématique les notions d'image inverse et d'image directe de cycles.
Dans le présent numéro, nous nous contenterons de définir ces notions dans certains cas particuliers utiles, et pour les cycles $1$-codimensionnels, ces définitions étant choisies de telle sorte qu'elles soient compatibles avec les définitions correspondantes pour les diviseurs \sref{IV.21.4} et \sref{IV.21.5}, compte tenu de l'homomorphisme $\operatorname{cyc}$ défini en \sref{IV.21.6}.

\begin{env}[21.10.1]
\label{IV.21.10.1}
